\section{Guenon and Tradition}

Tradition, in the Guenonian sense, has nothing to do with the past in itself. Just because it happened yesterday does not make it therefore traditional. And it is not a reaction against modernity, nor a nostalgia for the past. Its expression today may appear that way, but that is because a Traditional path either no longer exists in the West, or is accessible only to a few.

The method of Tradition has nothing at all to do with the study of comparative religion. One can be a so-called ``Traditionalist'' within the confines of one's own religion alone, without any detailed knowledge of any other religion whatsoever. As a matter of fact, that is the way it has been for most of history, recorded and unrecorded. So it is not a process of somehow ``distilling'' the core teachings of several religions to end up with something remarkable. I don't see what that even entails or how it could be accomplished, since it implies one can somehow have the inside (``esoteric'') without an outside (``exoteric''). Even worse, is syncretisim, that is, combining all the religions into some incompatible stew.

So that brings us to the next item: the relationship between the esoteric and exoteric sides of a religious tradition. In \emph{The Esoterism of Dante}, Guenon responds to the question of whether Dante was Christian or pagan. Guenon explains:

\begin{quotex}
we do not think that such a point of view is necessary, for true esoterism is something completely different from outward religions, and if it has some relationship with it, this can only be insofar as it finds a symbolic mode of expression in religious forms. Moreover, it matters little whether these forms be of this or that religion, since what is involved is the essential doctrinal unity concealed beneath their apparent diversity.
\end{quotex}


So this brings us to the crux of the matter, which few people seem to find palatable: ``in the past, initiates participated in all forms of worship, following the customs established in whatever country they happened to be.'' So let us make clear the necessary corollaries that follow from this.

Esoterism is neither pagan nor Christian. So whoever --- and I have met many --- believes that by rejecting Christianity and reverting to some pagan practice, whether Greek, Roman, or Germanic, he therefore is coming closer to Tradition, then he is quite mistaken. Guenon again:

\begin{quotex}
The ancient mysteries were not paganism, but were superimposed upon it. In the same way there were in the Middle Ages some organizations of an initiatic, and not religious, character, but which took Catholicism as their base.
\end{quotex}


Since Catholicism has been, at least until recently, the Traditional form of worship in Europe, then we would expect Traditionalists to find their home there. The idea that we can ``pick and choose'' our religion comes from a modernist mindset and is alien to the Traditionalist. Unfortunately, Catholicism has lost much of its Traditional character, leaving Traditionalists without solid ground in the West; reenacting the lost forms of ancient paganisms is an even worse option.


\flrightit{Posted on 2010-02-23 by Cologero}

